\chapter{计算机模拟}
计算机模拟在现代统计力学中发挥着重要作用。计算机模拟在科学领域的应用历史，与早期数字计算的历史相辅相成。参与这一领域的人们和地点主要集中在洛斯阿拉莫斯及其他美国国家实验室，这些地方在第二次世界大战结束后，首批数字计算机得以投入使用，供科学家们使用。在计算机模拟方法的开发初期，费米、乌拉姆、冯·诺依曼、泰勒、梅特罗波利斯、罗森布拉特等人均作出了重要贡献，他们同样参与了曼哈顿计划（梅特罗波利斯，1987年）。
统计力学中的计算机模拟主要分为两大类：蒙特卡洛（MC）法和分子动力学（MD）法，尽管二者之间还存在多种变体。这两种方法都通过数值演化材料的简单模型，依据一组微观状态，来计算可测量物理量的热力学平均值。计算机模拟为研究物理系统提供了一种手段，这种手段既可与实验相辅相成，又可与理论相互补充。以下是计算机模拟的若干优势：
计算机模拟能够帮助我们深入理解模型系统的平衡态与非平衡态行为，尤其适用于那些理论近似失效或尚未经过验证的参数范围。
计算机模拟为检验理论近似在特定模型系统中的适用范围提供了有效途径。
计算机模拟能够实现对物理过程的可视化，从而为复杂现象带来全新的洞见。
计算机模拟使我们能够对那些在实验中难以直接观察到的行为进行细致探究。
计算机模拟可用于研究基础物理过程，这些过程可以为理论发展提供重要指导。
计算机模拟能够用于构建自然界中并不存在的系统，以帮助我们更好地理解现有材料，并设计和制造全新的材料。
\section{引言与统计学}\label{ux5f15ux8a00ux4e0eux7edfux8ba1ux5b66}
虽然计算机模拟理论中某些关键环节应严格遵循，但计算机模拟的开发与应用很大程度上仍是一种艺术形式。针对任何特定问题，都有多种可能的模拟方法可供选择，而其中某些方案往往效果更佳。
在阐明重要物理性质方面，其作用往往优于其他方法。本章简要介绍平衡模拟，但计算机模拟也被广泛应用于动力学过程和非平衡过程的建模。要确定模型系统的平衡热力学平均值，需要生成一系列微态，并从模型的平衡态集合中进行选择。例如，可以通过MD模拟来求解牛顿运动方程，从而在相空间中生成状态的时间序列，随着系统不断探索哈密顿量的恒能超曲面。相比之下，对同一模型进行MC模拟时，可能会生成一系列状态，这些状态是通过在规范态集合的构型微态之间进行随机游走而选定的。这两种方法都是重要抽样法的典型实例：重要抽样法将计算工作重点放在生成能够代表平衡态集合的微态上，而非对整个相空间进行采样。正是这种效率上的巨大提升，使得统计力学模型的计算机模拟成为可能。无论采用哪种方法生成的状态序列，都可以用来估算平衡平均值。Allen 和 Tildesley（1990）、Binder 和 Heermann（2002）、Frenkel 和 Smit（2002）以及 Landau 和 Binder（2009）都对计算机模拟及其在统计物理学中的应用进行了更为详尽的探讨。
设 \(q\) 为系统的某一微态，\(A(q)\) 为与该微态相关的热力学可观测量。在MC模拟中，\(q\) 可能代表系统中所有粒子的位置；而在MD模拟中，\(q\) 可能代表所有粒子的位置和动量。可观测量 \(A(q)\) 可以代表势能、压力的维里贡献、配对相关函数等。用于启动模拟的初始微态通常并不一定代表构成平衡态集合的全部微态，但模拟的目标是使微态在平衡态集合的足够大子集内演化，从而使可观测量的平均值逐渐接近其平衡值。当模拟运行足够长，使系统接近平衡状态后，模拟会生成一组 \(M\) 个构型，即 \(\{q_j\}_{j=1}^M\)，这些构型来自平衡态集合中的微态集合，并为每个想要测量的热力学变量存储一组数值 \(\{A(q_j)\}_{j=1}^M\)。¹ 由于微态是从平衡态集合中选取的，因此 \(A\) 的平衡平均值可近似为一组数值 \(\{A(q_j)\}_{j=1}^M\) 的简单平均值。当然，模拟只能提供有限数量的状态序列，因此对结果不确定性的统计分析，是任何模拟中至关重要的组成部分。
变量 \(A\) 的平衡平均值由以下公式给出：
\begin{equation}\tag{16.1.1}\label{eq:16.1.1}
\langle A \rangle = \langle A \rangle _ { M } \pm \sigma _ { M }
\end{equation}
\({}^{1}\)或者，也可以将所有统计独立的微态配置集存储起来，以供后续分析使用。这种策略需要占用大量存储空间，但在生成统计独立配置的成本较高、且在日后需要利用已存储的配置来计算多个不同可观测量的平均值时，会非常有用。这种方法有时会被应用于大规模晶格量子色动力学模拟中。
其中，模拟平均值 \(\langle A\rangle_{M}\) 与不确定性 \(\sigma_{M}\) 分别由以下公式确定：
\begin{equation}\tag{16.1.2}\label{eq:16.1.2}
\langle A \rangle _ { M } = \frac { 1 } { M } \sum _ { j = 1 } ^{M} A ( q _ { j } ) ,
\end{equation}
\begin{equation}\tag{16.1.3}\label{eq:16.1.3}
\sigma _ { M } = \frac { \sqrt { \langle A ^{2} \rangle _ { M } - \langle A \rangle _ { M } ^{2} } } { \sqrt { M / ( 2 \tau + 1 ) } } \, ,
\end{equation}
\begin{equation}\tag{16.1.4}\label{eq:16.1.4}
\langle A ^{2} \rangle _ { M } - \langle A \rangle _ { M } ^{2} = \frac { 1 } { M } \sum _ { j = 1 } ^{M} [ A ( q _ { j } ) - \langle A \rangle _ { M } ] ^{2} .
\end{equation}
"相关时间"\(\tau\) 的定义如下：由于状态 \(q_j\) 是由模拟按顺序生成的，因此每个新的状态 \(q_{j+1}\) 都能确保与前一个状态 \(q_j\) 接近，因此序列中各状态 \(A(q_j)\) 的值具有高度相关性。随着"相关时间"\(\tau\) 的增加，\(A(q_j)\) 值之间的相关性会逐渐减弱，而\(\tau\) 可以通过相关函数 \(\phi_{AA}(t)\) 计算得出，即：
\begin{equation}\tag{16.1.5}\label{eq:16.1.5}
\phi _ { A A } ( t ) = \frac { \langle A ( t ) A ( 0 ) \rangle - \langle A ( t ) \rangle \, \langle A ( 0 ) \rangle } { \left\langle A ^{2} \right\rangle - \langle A \rangle ^{2} } ,
\end{equation}
\begin{equation}\tag{16.1.6}\label{eq:16.1.6}
\tau = \sum _ { t > 0 } \phi _ { A A } ( t ) .
\end{equation}
变量 \(t\) 表示有序序列中各配置对之间的分离程度。在分子动力学模拟中，\(\tau\) 代表系统在能量面上沿其轨迹移动足够远所需的时间，从而使得 \(A\) 的值不再存在相关性。蒙特卡洛模拟则以随机游走的方式探索平衡微态，因此 \(\tau\) 并不对应于物理时间，而是指完成统计独立的 \(A\) 值所需的蒙特卡洛采样次数的平均值。量 \(M/(2\tau + 1)\) 表示在 \(M\) 个数值的序列中，统计独立的配置数量。²
\(^{2}\)式\autoref{eq:16.1.3} 中的相关函数 \(\phi_{AA}(t)\) 可以通过 \(M\) 个配置的子序列来测量：
\begin{equation}\tag{16.1.7}\label{eq:16.1.7}
\langle A ( t ) A ( 0 ) \rangle \approx \frac { 1 } { M ^{\prime} } \sum _ { j = 1 } ^{M ^ { \prime} } A ( q _ { j + t } ) A ( q _ { j } ) ,
\end{equation}
并且
\begin{equation}\tag{16.1.8}\label{eq:16.1.8}
\langle A ( t ) \rangle \approx \frac { 1 } { M ^{\prime} } \sum _ { j = 1 } ^{M ^ { \prime} } A ( q _ { j + t } ) .
\end{equation}
根据定义，相关函数 \(\phi_{AA}(0)\) 的值为 1，而当 \(t \to \infty\) 时，相关性会逐渐衰减至零。一旦我们能够根据系统的平衡相关性知识，为某一特定系统合理地估算出相关时间 \(\tau\) 的上限，就可以在存储 \(A(q_j)\) 值之间，简单地跳过超过 \(\tau\) 个配置，从而确保序列中的数值如今大致具备统计独立性。
\section{蒙特卡洛模拟}\label{ux8499ux7279ux5361ux6d1bux6a21ux62df}
"蒙特卡罗方法"这一术语以摩纳哥的赌场命名，由尼古拉斯·梅特洛普斯（1987）提出——"这一说法并非与斯坦·乌拉姆有一位叔叔有关，因为这位叔叔总是向亲戚借钱，只因他`非去摩纳哥不可'。"在平衡统计力学中，蒙特卡罗模拟的目标是利用伪随机数，从平衡概率分布中抽取一组具有代表性的微态\{q\}。
\begin{equation}\tag{16.2.1}\label{eq:16.2.1}
P _ { \mathrm { e q } } ( q ) = \frac { \exp ( - \beta E ( q ) ) } { \sum _ { q ^{\prime} } \exp [ - \beta E ( q ^{\prime} ) ] } .
\end{equation}
这意味着，"我们不再随机选择构型，再用\(\exp(-E/kT)\)对它们进行加权，而是以\(\exp(-E/kT)\)的概率来选择构型，并对这些构型进行均匀加权"（梅特洛普斯等人，1953）。如果一次模拟能够实现这一点，那么就可以利用式\autoref{eq:16.1.2}中的简单平均值来计算热力学平均值。这种对状态进行重要性采样的方法，相较于常规的随机采样，在计算上具有巨大的优势。
以下算法旨在实现这样一个目标：从所有微态集合中随机选取微态\(q\)，其概率分布接近于平衡分布（1）（梅特洛普斯等人，1953；卡洛斯与惠特洛克，1986；艾伦与蒂尔德斯利，1990；弗伦克尔与斯密特，2002；宾德与赫尔曼，2002；兰道与宾德，2009）。假设一个微态集合拥有初始的概率分布\(P(q,0)\)，并让该分布按照离散的随机速率方程演化。
\begin{equation}\tag{16.2.2}\label{eq:16.2.2}
P ( q , t + 1 ) = P ( q , t ) + \sum _ { q ^{\prime} } P ( q ^{\prime} , t ) W ( q ^{\prime} \rightarrow q ) - P ( q , t ) \sum _ { q ^{\prime} } W ( q \rightarrow q ^{\prime} ) ,
\end{equation}
其中\(W(q \rightarrow q')\)表示从状态\(q\)到状态\(q'\)的跃迁速率。若跃迁速率满足平衡条件
\begin{equation}\tag{16.2.3}\label{eq:16.2.3}
\sum _ { q ^{\prime} } P _ { \mathrm { e q } } ( q ^{\prime} ) W ( q ^{\prime} \rightarrow q ) = P _ { \mathrm { e q } } ( q ) \sum _ { q ^{\prime} } W ( q \rightarrow q ^{\prime} ) ,
\end{equation}
并且，若式（2）中的随机过程能够在有限步数内，从每一个微态到达每一个其他微态，那么该集合的概率将逐渐逼近平衡分布：
\begin{equation}\tag{16.2.4}\label{eq:16.2.4}
\operatorname* { l i m } _ { t \to \infty } P ( q , t ) = P _ { \mathrm { e q } } ( q ) .
\end{equation}
在实践中，式（3）通常采用详细平衡条件来实现。
\begin{equation}\tag{16.2.5}\label{eq:16.2.5}
P _ { \mathrm { e q } } ( q ) W ( q \rightarrow q ^{\prime} ) = P _ { \mathrm { e q } } ( q ^{\prime} ) W ( q ^{\prime} \rightarrow q ) .
\end{equation}
要计算\(P_{\mathrm{eq}}(q)\)，需要对所有状态求和以确定配分函数，但\(P_{\mathrm{eq}}(q')/P_{\mathrm{eq}}(q)\)的比值仅取决于能量差\(\Delta E = E(q') - E(q)\)。
因此，跃迁速率之间存在如下关系：
\begin{equation}\tag{16.2.6}\label{eq:16.2.6}
W ( q \rightarrow q ^{\prime} ) = \exp ( - \beta \Delta E ) \, W ( q ^{\prime} \rightarrow q ) .
\end{equation}
这保证了，无论从何种起始构型出发，由这一随机过程生成的状态序列，最终会趋于等价——即通过在平衡集合的微态间进行随机游走来选择状态。这一方法可以借助计算机代码实现，正如梅特洛普斯等人（1953）首次提出的那样，利用跃迁速率来完成。
\begin{equation}\tag{16.2.7}\label{eq:16.2.7}
\begin{array}{rl} & { W ( q \to q ^{\prime} ) = 1  \mathrm { ~ i f ~ } \Delta E \leq 0 , } \\& { W ( q \to q ^{\prime} ) = \exp ( - \beta \Delta E )  \mathrm { ~ i f ~ } \Delta E > 0 . } \end{array}
\end{equation}
过渡率的其他选择是可能的，但这种以梅特罗波利斯命名的形式，是使用最广泛的其中一种。
\section{A 梅特罗波利斯蒙特卡洛算法}\label{a-ux6885ux7279ux7f57ux6ce2ux5229ux65afux8499ux7279ux5361ux6d1bux7b97ux6cd5}
梅特罗波利斯方法可以通过计算机程序实现，利用伪随机数生成器 rand() 来生成伪随机数，这些随机数在开区间 (0.0, 1.0) 上均匀分布；有关如何生成伪随机数的详细讨论，请参见附录 I。首先，通过从模型的所有微态集合中选择一个起始状态 \(q_0\) 来初始化系统。如果 \(q_0\) 不属于平衡态集合中的典型状态，则会更有助于系统达到平衡。这样可以减少系统达到平衡所需的步数。例如，一个无序的液态状态并不适合作为晶体固体模拟的起始点。
梅特罗波利斯算法由以下步骤定义：
生成一个随机的试验态 \(q_{\mathrm{trial}}\)，该态与系统的当前态 \(q_{j}\) "相近"。这里的"相近"意味着，试验态应与当前态几乎完全相同，仅在少数随机变化上有所差异，通常是对单个粒子或自旋进行微小调整。例如，可以通过随机将一个粒子移动到邻近位置，来创建粒子模拟的试验态。
\begin{equation}\tag{16.2.8}\label{eq:16.2.8}
x _ { i } ^{\mathrm { t r i a l} } = x _ { i } + \Delta x ( \mathrm { r a n d } ( ) - 0 . 5 ) ,
\end{equation}
再调用两次 rand()，以生成 \(y_{i}^{\mathrm{trial}}\) 和 \(z_{i}^{\mathrm{trial}}\)。对于自旋系统而言，试验态通常涉及自旋翻转，或是对单个自旋进行随机旋转。³
³存在用于自旋系统的蒙特卡洛算法，这些算法能够对大型相关簇中的自旋进行翻转，而非逐个翻转单个自旋；请参阅 Swendsen 和 Wang（1987）以及 Wolff（1989）。这些方法在模拟某些特定模型时非常有效。此外，由于每次翻转的接受与否与其他已翻转自旋无关，因此也可以尝试一次性对非相互作用子格中的所有自旋进行翻转。例如，在某些计算机架构或编程环境中，对自旋模型进行棋盘式更新——即自旋仅与正方形晶格的最近邻自旋相互作用——可能会更加高效。
计算试验态与前一态之间的能量变化，即 \(\Delta E = E(q_{\text{trial}}) - E(q_j)\)。若 \(\Delta E \leq 0\)，则接受该试验态，即设置 \(q_{j+1} = q_{\text{trial}}\)。若 \(\Delta E > 0\)，则以概率 \(\exp(-\beta \Delta E)\) 接受该试验态。实现这一过程需要额外调用伪随机数生成器。若 \(\text{rand}(.) < \exp(-\beta \Delta E)\)，则接受该试验态。如果相互作用为短程相互作用，那么能量变化的计算仅涉及与附近少数几个粒子或自旋之间的相互作用。若试验态在某种方面存在非法性，即它并非所有构型集合中的合法状态，则应拒绝该状态。这相当于将能量变化设为 \(+\infty\)。若试验态因任何原因被拒绝，则将系统的新的状态设置为前一状态 \(q_{j+1} = q_j\)，即保持状态原值 \(q_j\)，丢弃该试验态，并继续下一步。
对系统中的每个粒子或自旋执行一次步骤 1 和步骤 2。通常采用随机方式完成此操作，以确保详细平衡。⁴ 步骤 1 至步骤 3 定义了一次蒙特卡洛扫描。
重复步骤 1 至步骤 3，进行 \(M_{\text{eq}}\) 次蒙特卡洛扫描，以使系统达到平衡。事先很难明确选择 \(M_{\text{eq}}\) 的合适值。至少，在模拟过程中所研究的所有量 \(A(q)\) 在平衡结束时应不再出现明显的单调漂移。然而，这并不能保证系统已真正达到平衡，因为系统很可能仍被困在一种长寿命的亚稳态附近。
在进行 \(M\) 次蒙特卡洛扫描时，重复步骤 1 至步骤 3，并持续记录所有想要测量的热力学变量，即 \(\{A(q_j)\}_{j=1}^M\)。利用式\autoref{eq:16.1.1} 和 式\autoref{eq:16.1.2}，确定系统的平衡平均值及不确定性。
若要根据不同的参数集（如温度、密度等）确定平均值，只需将参数略微调整，然后重复步骤 1 至步骤 5，包括平衡步骤 4.5。将上一次运行的最后一个构型作为下一次运行的首个构型，往往能够显著缩短平衡时间。图~16.1 展示了在 \(128 \times 128\) 方形晶格上，对二维伊辛模型的比热进行的蒙特卡洛计算结果，与第 13.4.A 节中给出的精确解进行了对比。
⁴ 有时会采用违反详细平衡的顺序更新方法及其他更新方式以提高效率，但必须格外谨慎，确保在平均意义上维持详细平衡。
在某些情况下，可以利用直方图重采样方法来减少需要模拟的温度和场的数量（Ferrenberg 和 Swendsen，1988）。例如，如果在耦合参数 \(K\) 和磁场 \(h\) 下进行的自旋模拟，得到了一个能够采样联合能量-磁化分布 \(P_{K,h}(E,M)\) 的直方图，那么在邻近温度和场下的分布可由以下公式给出：
\begin{equation}\tag{16.2.9}\label{eq:16.2.9}
P _ { K + \Delta K , h + \Delta h } ( E , M ) = \frac { P _ { K , h } ( E , M ) e ^{\Delta K E + \Delta h M} } { \sum _ { E , M } P _ { K , h } ( E , M ) e ^{\Delta K E + \Delta h M} } .
\end{equation}
目前广泛采用的其他方法包括：并行退火、多标度蒙特卡洛以及"宽直方图"方法。这些方法在研究具有强第一类相变的系统时尤为有效。蒙特卡洛重整化群方法在研究临界点时具有极强的效能。有关综述，请参阅 Landau 和 Binder（2009）。
其中 \(F_{i}\) 是由 \(N\) 个粒子的势能函数 \(U\) 对粒子 \(i\) 所施加的力。分子动力学模拟通过离散的时间步 \(\Delta t\) 将系统向前推进。最
在此背景下，常用的积分方法源自 Verlet（1967）：
\begin{equation}\tag{16.2.10}\label{eq:16.2.10}
\begin{array}{r} { r _ { i } ( t + \Delta t ) = 2 r _ { i } ( t ) - r _ { i } ( t - \Delta t ) + \frac { ( \Delta t ) ^{2} } { m _ { i } } F _ { i } ( t ) . } \end{array}
\end{equation}
这一方法等同于"跳蛙"算法和速度Verlet算法，它们同时更新粒子的位置和速度；详情请参见Frenkel和Smit（2002）。Verlet方法能够保留哈密顿运动方程的时间反演对称性，且每步的误差阶为\((\Delta t)^4\)，同时只需对每个粒子进行一次力的计算——而这通常是模拟中耗时最长的环节。最重要的是，Verlet算法具有辛结构，因此其积分过程相当于对"附近"的"幽灵"哈密顿量求解精确解，从而实现了良好的长期稳定性，并有效保护了系统的能量特性。
在模拟开始时，系统会以某种初始微态启动，其中所有粒子的位置和速度均已明确设定；随后，积分算法会按照时间顺序向前推进粒子的位置和速度。用于近似计算粒子速度和能量的最简单形式包括：
\begin{equation}\tag{16.2.11}\label{eq:16.2.11}
\pmb { v } _ { i } ( t ) = \frac { \pmb { r } _ { i } ( t + \Delta t ) - \pmb { r } _ { i } ( t - \Delta t ) } { 2 \Delta t } ,
\end{equation}
\begin{equation}\tag{16.2.12}\label{eq:16.2.12}
E = \sum _ { i = 1 } ^{N} \frac { m _ { i } } { 2 } \left( \frac { \pmb { r } _ { i } ( t + \Delta t ) - \pmb { r } _ { i } ( t - \Delta t ) } { 2 \Delta t } \right) ^{2} + U [ \pmb { r } _ { 1 } ( t ) , \pmb { r } _ { 2 } ( t ) , . . , \pmb { r } _ { N } ( t ) ] .
\end{equation}
时间步长\(\Delta t\)应选择得足够小，以与哈密顿体系中最短的基本时间尺度相匹配，但又不能过小以至于影响程序的运行效率。数值积分近似地描述了微正则系综中的一个成员，在相空间中沿恒能超曲面运动。要正确计算平衡平均值，关键在于哈密顿量必须是遍历的；⁶ 这一特性使得系统能够采样恒能超曲面中的所有区域，因此MD时间平均值等同于微正则系综的平均值。如果在模拟过程中，总能量偏离了预先设定的某个阈值，那么可以对所有速度进行重新缩放，使总能量恢复到初始值。或者，也可以使用恒温器将温度维持在所需目标值；详情请参见Frenkel和Smit（2002）。
\(^{6}\)由于分子动力学模拟会生成模型系统的时序演化，因此我们需要确保系统具备遍历性——也就是说，时间平均与集合平均应相等。例如，谐振子系统并不具备遍历性。一个在\(d\)维空间中包含\(N\)个粒子的系统，其相空间维度为\(6N\)，因此常能能超曲面的维度为\(6N-1\)。对于\(N\)个耦合振子的正常模解法，有\(3N\)个运动常数，因此系统仅探索\(3N\)维的超曲面。即便将粒子间的相互作用设计成非谐振子，这一问题依然无法消除——正如费米、帕斯塔和乌拉姆于1955年首次揭示的那样，并且在理论层面被科尔莫戈洛夫（1954年）、阿诺德（1963年）以及莫泽尔（1962年）进一步深入探讨。平衡态系统的分子动力学模拟假定系统具有遍历性。目前仅有少数系统被证明是遍历性的；幸运的是，大多数具有真实配对势的系统，在两维或更多维空间中似乎都表现出遍历性。鉴于此，从这一视角来看，一维系统的分子动力学模拟应被视为有待商榷。
对于氖、氩等单原子流体，常用的一种配对相互作用是莱纳德-琼斯相互作用。
\begin{equation}\tag{16.2.13}\label{eq:16.2.13}
u ( r ) = 4 \varepsilon \left( \left( \frac { D } { r } \right) ^{12} - \left( \frac { D } { r } \right) ^{6} \right) ,
\end{equation}
其中，\(D\)为分子直径，\(\varepsilon\)为吸引势阱的深度。莱纳德-琼斯势在长距离处呈吸引力，且随着距离的增加而按\(1/r^6\)衰减，以模拟范德华力；而在短距离处，它则按\(1/r^{12}\)发散，以模拟泡利排斥力，从而防止电子波函数发生重叠。模拟最好采用无量纲参数进行。对于具有莱纳德-琼斯相互作用的流体，所有长度均可用\(D\)作为单位进行测量，所有能量（包括\(kT\)）均以\(\varepsilon\)为单位，所有力均以\(\varepsilon/D\)为单位，所有压力均以\(\varepsilon/D^3\)为单位，所有时间均以\(\sqrt{mD^2/\varepsilon}\)为单位，以此类推。随后，模拟只需针对缩减温度\(kT/\varepsilon\)、缩减密度\(nD^3\)等单一数值开展，同时以缩减后的单位来测量各种可观测物理量。之后，便可利用实验所测得的\(D\)、\(m\)和\(\varepsilon\)等参数，对模拟结果与实验数据进行对比分析。在无量纲单位下，一对粒子之间的莱纳德-琼斯力为
\begin{equation}\tag{16.2.14}\label{eq:16.2.14}
F = \pm \frac { r } { r ^{2} } \left( \frac { 48 } { r ^{12} } - \frac { 24 } { r ^{6} } \right) .
\end{equation}
牛顿第三定律可用于将力的计算次数减少一半。莱纳德-琼斯模型最早由伍德和帕克（1957年）在蒙特卡罗模拟中进行研究，随后又由拉赫曼（1964年）和弗雷特（1967年）在分子动力学模拟中得到深入探讨。
\section{A 分子动力学算法}\label{a-ux5206ux5b50ux52a8ux529bux5b66ux7b97ux6cd5}
首先，通过设定所有粒子的初始位置和速度来启动系统。通常，初始速度是通过从麦克斯韦分布中选取每个粒子速度矢量的各个分量来确定的。在归一化单位下，这一过程可以表示为：
\begin{equation}\tag{16.3.1}\label{eq:16.3.1}
P _ { \mathrm { M a x w e l l } } ( v _ { x } ( 0 ) ) = \frac { 1 } { \sqrt { 2 \pi \, T } } \exp \left( - \frac { v _ { x } ^{2} ( 0 ) } { 2 \, T } \right) ;
\end{equation}
\(^{7}\)例如，适用于氩气的莱纳德-琼斯参数分别为\(\varepsilon/k=119.8\)K和\(D=0.3405\)nm（Levelt，1960；Rowley、Nicholson和Parsonage，1975）。为了减少每次时间步需要考虑的相互作用数量，相互作用势几乎总是会在分子间的有限距离处被截断。对于莱纳德-琼斯相互作用，最常见的方式是将势值设为\(r_{\mathrm{max}}=2.5D\)。如果将势能设置为在距离大于\(r_{\mathrm{max}}\)时为零，那么这将在势能图中留下一个微小的不连续点。为了消除这一问题，通常会将势能向上移动\(-u(2.5D)\simeq0.0163\varepsilon\)，使势能在\(r_{\mathrm{max}}\)处为零。这样，就可以在蒙特卡洛模拟与分子动力学模拟之间进行直接比较。若未对势能进行上述调整，则无法直接比较蒙特卡洛模拟与分子动力学模拟的结果，因为势能的不连续性会导致一种狄拉克函数形式的力，该力会影响分子动力学模拟中的运动，却不会影响蒙特卡洛模拟中的构型。在比较蒙特卡洛模拟与分子动力学模拟结果与实验数据时，必须同时考虑由这种偏移所引起的扰动，以及配对势能中缺失的尾部效应。
请参阅附录I，了解如何使用统一的伪随机数生成器从高斯分布中进行取样。随后，可利用初始速度在第一次时间步之后设定粒子的位置，即
\begin{equation}\tag{16.3.2}\label{eq:16.3.2}
\begin{array}{r} { \pmb { r } _ { i } ( \Delta t ) = \pmb { r } _ { i } ( 0 ) + \pmb { v } _ { i } ( 0 ) \Delta t + \frac { 1 } { 2 } \frac { ( \Delta t ) ^{2} } { m _ { i } } F _ { i } ( 0 ) . } \end{array}
\end{equation}
接下来，利用公式（2）将系统按照\(\tau_{\mathrm{eq}} / \Delta t\)的时间步向前推进。平衡时间\(\tau_{\mathrm{eq}}\)应选择得足够大，以确保系统能够达到平衡状态；具体细节请参见第16.2节中的蒙特卡洛讨论。通常会使用恒温器来使系统演化至所需温度的状态；详情请参考Frenkel和Smit（2002）。
现在，利用公式（2）将系统按照\(\tau_{\mathrm{avg}} / \Delta t\)的时间步向前推进，同时持续记录所有您希望测量的热力学变量\(\left\{A(q_{j})\right\}\)。最后，利用公式\autoref{eq:16.1.2}来计算系统的平衡平均值及不确定性。
若要根据不同的参数集（如温度、密度等）计算平均值，请对参数进行小幅调整，并重复步骤1和步骤2。将前一次运行的最后一个配置作为新运行的起始配置，往往能够缩短平衡时间。
\section{粒子模拟}\label{ux7c92ux5b50ux6a21ux62df}
通过在体积为 \(V\) 的周期性盒中放置 \(N\) 个粒子，并利用对偶势函数进行相互作用，可以采用蒙特卡罗模拟和分子动力学模拟来对流体进行建模。汉森与麦克唐纳（1986年）、艾伦与蒂尔德斯利（1990年）以及弗伦克尔与斯密特（2002年）对这一课题进行了详尽的综述。在蒙特卡罗模拟中，只需计算试探性运动所引起的能量变化；而在分子动力学模拟中，则只需计算粒子间最接近的周期性副本彼此距离在截断距离范围内的相互作用力。蒙特卡罗模拟通常采样规范态，因此能够控制温度和密度，并测量能量、压力等物理量；而分子动力学模拟通常采样微正态态，因此能够控制能量和密度，并测量温度、压力等物理量。根据平分定理，温度可表示为
\(d\) 维系统中的 \(T\)
\begin{equation}\tag{16.4.1}\label{eq:16.4.1}
k T = \frac { 1 } { N d } \left\langle \sum _ { i = 1 } ^{N} m _ { i } v _ { i } ^{2} \right\rangle .
\end{equation}
式（3.7.15） 和 式（10.7.11） 中的维里状态方程可用于确定任意一种模拟方式下的压力 \(P\)：
\begin{equation}\tag{16.4.2}\label{eq:16.4.2}
\frac { P } { n k T } = 1 + \frac { 1 } { N d k T } \left\langle \sum _ { i < j } F ( r _ { i j } ) \cdot r _ { i j } \right\rangle = 1 - \frac { n } { 2 d k T } \int \frac { d u } { d r } r g ( r ) d r .
\end{equation}
如第 10.7 节所讨论的那样，流体的配对关联函数可用于测量多种热力学性质，包括压力、等温压缩率 \(\kappa_{T}\) 以及散射结构因子 \(S(k)\)；参见式（10.7.18） 至 (10.7.21)。配对关联函数 \(g(r)\) 是由在模拟过程中定期收集所有粒子对之间的距离直方图得到的，同时需考虑周期性边界条件，并将直方图按比例缩放，其比例系数等于半径为 \(r\)、厚度为 \(\Delta r\) 的各层壳体的总体积。在三维空间中，配对关联函数可表示为
\begin{equation}\tag{16.4.3}\label{eq:16.4.3}
g ( r ) = \frac { 2 V } { N ^{2} \left( \frac { 4 \pi } { 3 } \right) \left[ ( r + \Delta r ) ^{3} - r ^{3} \right] } \left\langle \sum _ { i < j } ^{N} \Delta _ { \Delta r } ( r _ { i j } - r ) \right\rangle ,
\end{equation}
其中，阶跃函数 \(\Delta_{\Delta r}(\xi)\) 在 \(0 < \xi < \Delta r\) 时取值为 1，否则为 0。该表达式表示直方图中每个区间内事件数量与相同密度的理想气体所预期的平均事件数量之比。
\section{A 硬球的模拟}\label{a-ux786cux7403ux7684ux6a21ux62df}
硬球系统已在蒙特卡罗模拟和分子动力学模拟中得到了广泛研究，也是最早采用这两种方法开展研究的模型（Metropolis 等人，1953 年；Adler 与 Wainwright，1957 年、1959 年）。硬球的对偶势为
\begin{equation}\tag{16.4.4}\label{eq:16.4.4}
u ( r ) = \left\{ \begin{array}{ll} { 0 } & { \mathrm { f o r } \ r > D , } \\{ \infty } & { \mathrm { f o r } \ r \leq D . } \end{array} \right.
\end{equation}
对于本模型所采样的空间构型而言，温度是一个无关参数，因为该对势并不具有有限的能量尺度。要全面探索相图，只需改变缩减后的粒子数密度 \(nD^{d}\) 即可。所有热力学性质要么与温度无关，要么以一种看似简单的形式随温度变化。例如，对于硬球系统，其缩放后压力 \(P/nkT\) 仅取决于缩减后的粒子数密度 \(nD^{d}\)。硬球的密度通常表示为
以填充因子 \(\eta\) 为变量，即系统中实际被粒子占据的体积占比。在三维空间中，体积分数可表示为 \(\eta=\pi nD^{3}/6\)。由于对势存在奇点，无法使用维里方程（2）来计算压力，但可以通过硬球的维里状态方程来确定压力，即（10.7.12）。
用于硬球的蒙特卡罗代码相对简单，因为在试移动过程中，能量变化要么为零，要么为无穷大。若粒子的试位移距离任意其他粒子的距离小于或等于 \(D\)，则该试位移将被拒绝；反之，则予以接受。这是首次在计算机模拟中研究的统计物理模型（Metropolis 等人，1953 年）。
要实现硬球的分子动力学模拟，需要采用与标准分子动力学不同的方法。由于势函数不连续，有限差分积分法在此处行不通。相反，可以利用运动方程的精确解。每个粒子在恒定速度下沿直线运动，除了在两粒子发生碰撞的瞬间——即当它们之间的距离为 \(D\) 时。由于势函数具有奇异性，碰撞过程可以被唯一地按时间顺序排列。每次碰撞既守恒动能，也守恒动量，因此在每次碰撞后，粒子的速度可通过分析碰撞时刻两粒子中心之间的速度和位移矢量来确定。两碰撞粒子速度的变化为
\begin{equation}\tag{16.4.5}\label{eq:16.4.5}
\Delta \pmb { v } _ { i } = - \left. \Delta \pmb { v } _ { j } = \frac { - \left( \pmb { r } _ { i j } \cdot \pmb { v } _ { i j } \right) \pmb { r } _ { i j } } { D ^{2} } \right| _ { | \pmb { r } _ { i j } | = D } ,
\end{equation}
其中 \(r_{ij}\) 和 \(v_{ij}\) 分别代表两粒子的相对位置和相对速度。模拟通过从一次碰撞到下一次碰撞，按时间顺序推动粒子运动，并根据式（5）改变参与碰撞的粒子对的速度。这是统计物理领域首次实现的分子动力学方法应用（Alder 和 Wainwright，1957 年、1959 年）。
图~16.2展示了流体相中的配对关联函数和结构因子，图~16.3则给出了硬球模型的相图。在低密度下，平衡相为短程有序流体；而在高密度下，该模型的平衡相则为长程有序的面心立方固体。值得注意的是，一个模型要具备晶体相，并不要求存在吸引相互作用。然而，要形成液---汽共存线及临界点，则必须存在吸引相互作用。正因如此，硬球模型的低密度相通常被称为流体相，而非液体相，因为该模型并不具备液---汽共存线。液---固共存线处的液体体积分数和固体体积分数分别为\(\eta_{l} \simeq 0.491 \pm 0.002\)和\(\eta_{s} \simeq 0.543 \pm 0.002\)。液---固共存压力可表示为\(P_{ls}^{*} = P_{ls}D^{3}/kT \simeq 11.55 \pm 0.11\)；详情请参见Speedy（1997）。在\(\eta < \eta_{l}\)的低密度流体相中，缩减压力可由卡纳汉---斯塔林状态方程（10.3.25）准确拟合。
\begin{equation}\tag{16.4.6}\label{eq:16.4.6}
\frac { P } { n k T } = \frac { 1 + \eta + \eta ^{2} - \eta ^{3} } { ( 1 - \eta ) ^{3} } ,
\end{equation}
其中\(\eta^{*}(= \eta/\eta_{cp})\)为实际堆积密度与最大紧密堆积密度\(\eta_{cp}\)之比，即\(\pi\sqrt{2}/6\simeq0.7405\)（Speedy，1997；Frenkel和Smit，2002）。当密度接近紧密堆积密度时，固体相中的压力会趋于发散。该模型
对于密度介于 \(\eta_{l}\) 与随机紧密堆积体积分数 \(\eta_{rcp} \simeq 0.644 \pm 0.005\) 之间的体系，也呈现出一种亚稳态的无序相；参见 Rintoul 和 Torquato（1996）。有关硬球模型研究的综述，请参阅 Mulero 等人（2008）。
\section{计算机模拟的局限性}\label{ux8ba1ux7b97ux673aux6a21ux62dfux7684ux5c40ux9650ux6027}
计算机模拟在统计物理领域得到了广泛应用，并在我们对众多物理系统的理解中发挥了重要作用。模拟能够补充理论和实验，并为那些难以进行精确或近似理论分析的系统研究提供了诸多优势。然而，重要的是要充分认识到这一技术本身的固有局限性。
计算机模拟必然涉及有限数量的自由度，通常为数百至数千个粒子或自旋。这一规模远远不足以展现热力学大系统中所发生的多种行为。例如，在一个三维立方体盒中模拟由1000个粒子组成的致密系统，每个线性尺寸上仅约有10个粒子，因此，超过大约五个粒子直径范围内的相关性会受到周期性边界条件的影响。从此类研究中准确提取热力学行为往往需要对不同尺寸的系统序列进行模型的有限尺度标度行为分析。
计算机模拟必然受到模拟时间尺度的限制。典型的分子时间尺度约为 \(t_0 \approx a/\nu\)，其中 \(a\) 是微观尺度，\(\nu\) 是分子速度。在室温附近的原子尺度下，\(a \approx 0.1\) 纳米，\(\nu \approx 100\) 米/秒，由此可得 \(t_0 \approx 10^{-12}\) 秒。在 MD 模拟中，每次时间步长都会使系统在时间上向前移动一个距离 \(\Delta t\)，而为了能够准确地求解运动方程，该时间步长必须远小于 \(t_0\)——例如，\(\Delta t \approx 10^{-14}\) 秒。若一次模拟将系统向前推进 \(10^6\) 个时间步长，则所采样的物理时间仅为 10 纳秒，而这可能不足以覆盖问题中许多重要的时间尺度。当系统本身具有极其缓慢的时间尺度时，这一点尤为突出，例如在二阶相变附近出现的临界减速现象，以及在一阶相变附近出现的滞后现象。蒙特卡洛模拟同样面临类似的问题：模拟必须运行足够长的时间，以便模型能够探索相空间中足够大的区域，从而捕捉到平衡态行为。在有利的情况下，可以通过诸如粗粒化、簇更新方法、并行退火、多可容蒙特卡洛方法、蒙特卡洛重整化群等特殊模拟方法来缓解这一问题；详见脚注 5。
为了提高计算效率，粒子间的相互作用通常被高度简化，且相互作用范围往往被截断。对于长程相互作用（如格鲁尔-偶极范德华力等），需要进行重整化或采用摄动理论来加以处理。——为了提高计算效率，粒子间的相互作用通常被高度简化，且相互作用范围往往被截断。对于长程相互作用（如库仑相互作用、偶极相互作用、范德华力等），需要通过重整化或采用摄动理论来加以处理，以尝试充分考虑其影响。
MC和MD模拟并不能直接测量系统可用的微观状态数，因此无法像计算其他可观测物理量那样直接求取熵或自由能。例如，若要确定自由能以定位相变，可通过对系统进行热力学积分，将其推导至具有理论已知自由能的状态；参见Frenkel和Smit（2002）。
MD模拟依赖于哈密顿量的遍历性，因此一维模型、从机械平衡中微弱受扰的模型，以及其他近乎可积的模型，可能会被困在低维轨道中，而无法完全探索哈密顿量的恒能超曲面；参见脚注6。
所有伪随机数生成器在其序列中都会产生一定程度的相关性。即使伪随机数序列中的细微相关性，也可能在蒙特卡洛模拟中导致错误的结果。不同类型的生成器存在不同的弱点，因此如果切换到基于不同算法的生成器，有时可以解决由特定生成器产生的相关性所引发的问题。在将生成器用于模拟之前进行测试，始终是一个明智的选择；参见附录I。
确认MC和MD模拟代码的有效性至关重要。这一过程与验证理论计算的方式大不相同。一些代码评估与验证流程包括：首先在已知性质的小规模系统上对代码进行测试；尽可能在有精确解的模型或已在文献中广泛研究过的模型上对代码进行测试；考察结果随系统尺寸和运行时间的变化；并在新增交互项或代码模块时，进行仔细的重新测试。在这方面，Frenkel和Smit（2002）、Parker（2008）以及Landau和Binder（2009）都提供了诸多有效的策略清单。
\section{问题}\label{ux95eeux9898_chap16}
16.1. 编写一段代码，用于测试均匀伪随机数生成器。若手头没有现成的生成器，可依据附录I中L'Ecuyer推荐的生成器进行编写。请执行以下测试：平均值 \(\langle x \rangle = 1/2\)，方差 \(\langle x^2 \rangle - \langle x \rangle^2 = 1/12\)，以及双相关性检验 \(\langle x_{i+k}x_i \rangle = 1/4\)（对于 \(k \neq 0\)）。在二维单位正方形上生成一组数字对的直方图，并检验其分布是否在统计学上呈均匀分布。
16.2. 编写一段代码，用于测试高斯伪随机数生成器。若手头没有现成的生成器，可依据附录I中的Box-Muller算法进行编写。请执行以下测试：平均值 \(\langle x \rangle = 0\)，方差 \(\langle x^2 \rangle = 1\)，以及双相关性检验 \(\langle x_{i+k}x_i \rangle = 0\)（对于 \(k \neq 0\)）。在二维空间中生成一组数字对的直方图，并检验其分布是否在统计学上呈高斯分布。
16.3. 定义一个具有相关性的随机数序列
\begin{equation}\tag{16.5.1}\label{eq:16.5.1}
s _ { k } = \alpha s _ { k - 1 } + ( 1 - \alpha ) r _ { k } ,
\end{equation}
其中 \(r_{k}\) 是一个方差为1、彼此不相关的高斯伪随机数，而 \(0 < \alpha < 1\) 则定义了相关性的范围。证明该序列服从高斯分布，且其均值为零。
均值。根据 \(\alpha\) 计算方差，并将所得结果与式\autoref{eq:16.1.2} 进行比较。编写一段代码，用于计算相关函数\autoref{eq:16.1.3}。绘制你所测量的相关函数的图形，并将其与精确的相关函数进行对比。
16.4. 编写一段蒙特卡洛代码，用于模拟一个由 \(N\) 个直径为 \(D\) 的硬球组成的系统，这些硬球位于长度为 \(L\) 的一维环形结构中，并采用周期性边界条件。计算两两之间的相关函数，并将其与式\autoref{eq:13.1.6} 和式\autoref{eq:13.1.7} 进行比较。系统的压力由公式 \(P/nkT = 1 + nDg(D^+)\) 给出；参见式\autoref{eq:10.7.12}。将你的压力与基于精确构型配分函数 \(Z_N = L(L - ND)^{N-1}/N!\) 所得的压力进行比较；参见式\autoref{eq:13.1.2}。
16.5. 编写一段蒙特卡洛代码，用于模拟一个由 \(N\) 个硬球组成的流体系统，该流体系统被置于一个长宽均为 \(L \times L\) 的二维方形盒中，并在每个方向上均采用周期性边界条件。计算两两之间的相关函数，并利用式\autoref{eq:10.7.12} 确定缩放后压力，即 \(P/nkT = 1 + 2\eta g(D^{+})\)。将此压力与近似形式 \(P/nkT = (1 + \eta/8)/(1 - \eta)^{2}\) 进行比较。
16.6. 为一个由 \(N\) 个刚性球组成的二维 \(L \times L\) 方形盒中的流体编写蒙特卡洛代码，并在算法中加入单粒子引力项 \(\sum_{i=1}^{N} mgy_{i}\)，即以 \(\exp(-\beta m g \Delta y)\) 的概率接受那些在物理上符合约束条件的构型。您需要在上下两面墙采用硬壁边界条件。求解平均密度随盒子垂直位置的变化关系。
16.7. 为一个由 \(N\) 个莱纳德-琼斯粒子组成的二维 \(L \times L\) 方形盒编写分子动力学代码。在每个方向上均采用周期性边界条件。利用维里方程\autoref{eq:16.4.2}计算缩放压力。计算并绘制系统的配对关联函数。
16.8. 为一个由 \(N\) 个莱纳德-琼斯粒子组成的二维 \(L \times L\) 方形盒编写分子动力学代码，并在能量中加入单粒子引力项：\(\sum_{i=1}^{N} mgy_{i}\)。在 \(x\) 方向采用周期性边界条件，但在上下两面墙则采用斥力 WCA（Weeks、Chandler 和 Andersen，1971）势。WCA 势是莱纳德-琼斯势中对于 \(r/D < (2)^{1/6}\) 的斥力部分，且该势会向上偏移 \(\varepsilon\)。证明：每个粒子的平均动能与盒子高度 \(y\) 无关，但平均缩放密度 \(nD^{2}\) 取决于盒子的垂直位置。
16.9. 编写 MC 代码，模拟一维伊辛模型，其在长度为 \(L\) 的周期性晶格上运行。计算该模型的内能和比热，并将其与方程（13.2.15）和（13.2.16）进行比较。计算相关函数 \(G(n)=\langle s_{i+n}s_{i}\rangle\)，并与方程（13.2.32）进行对比。
16.10. 编写 MC 代码，模拟二维邻近邻居伊辛模型，其在零场下运行于一个周期性 \(L \times L\) 晶格上。计算系统的内能和比热随温度的变化，并将其与第 13.4.A 节中的精确结果进行比较。有关二维伊辛模型在不同晶格尺寸下的精确结果，请访问 wwwELSEvierdirect.com。
16.11. 编写 MC 代码，模拟二维邻近邻居伊辛模型，其在零场下运行于一个周期性 \(L \times L\) 晶格上。计算能量分布 \(P(E)\)，涵盖多个温度范围，包括临界点。利用该分布，计算内能和比热随温度的变化。有关二维伊辛模型在不同晶格尺寸下的精确结果，请访问 www.elsevierdirect.com。
16.12. 编写 MC 代码，模拟一维 XY 模型。计算内能、比热、等温磁化率以及配对关联函数，并将您的结果与第 13.2 节中 \(n=2\) 情况的解析结果进行比较。
16.13. 编写一个MC代码来模拟二维XY模型。计算内能、比热、等温磁化率以及配对关联函数，并将你的结果与第13.7节中给出的理论结果进行比较。



